\section{Reproducibility and repository integration}
\label{sec:reproduce}

\subsection{What the archive contains}

The archive includes the complete article source, generated tables and
scientific figures, and the compiled PDF. A curated \code{source/} tree
contains the maintained Python implementation, its tests and fixtures, and
the historical reference modules required by those tests. Historical result
archives unrelated to this continuation are omitted. The distributed source
is pinned to \sourcepin{} except for the explicit delta in
\code{source_manifest.json}.

The \code{patches/} directory contains a Git-applicable patch with two
modified files and fourteen new files. A separate \code{baseline/} directory
retains the two modified files exactly as they appeared at the pinned
revision. It lets the scalar benchmark load its historical arm without
network access or a Git checkout. Independent validation applied the patch
in a fresh directory, checked all sixteen affected outputs byte for byte,
and compared every one of the 362 distributed source files with the working
tree. The 346 unmodified source files match the pinned Git blobs exactly.

The \code{experiments/} directory contains the arithmetic and geometric
reference checks, including rational geometry and optional Regina
comparisons. \code{results/} preserves the recorded measurements and exact
comparison counts. \code{validation/} preserves the baseline failure,
final 790-test success, targeted checks, patch audit, and artifact checks.
Standalone copies of the small arithmetic and extraction modules are kept
only to make those experiment scripts independent of package installation;
their bytes match the corresponding integrated modules.

\subsection{Commands for independent reproduction}

The new mathematical kernels require only the Python standard library.
Regina is optional for the independent native audit; Matplotlib is optional
for regenerating the figures. The recorded versions are listed in
\code{requirements-audit.txt}. LaTeX compilation uses ordinary
\code{pdflatex} and \code{latexmk}, with the packages declared in the main
source file.

\noindent\begin{minipage}{\linewidth}
From the extracted archive root:

\begin{lstlisting}
python scripts/check_hashes.py
python scripts/reproduce.py quick

# Optional native audit and figure dependencies:
python -m pip install -r requirements-audit.txt
python scripts/reproduce.py full

# Paired and binary-size benchmarks:
python scripts/reproduce.py benchmarks

# Rebuild the article:
cd article
latexmk -pdf -interaction=nonstopmode -halt-on-error \
  unknot_affine_families.tex
\end{lstlisting}
\end{minipage}

The reproduction driver writes new logs and measurements to
\code{reproduction/}; it preserves the recorded evidence. ``Quick'' runs
all 65 methods in the five new or restored test modules, the rational
face-order audit, the certificate examples, and the peripheral reduction
check. ``Full'' runs the complete integrated suite and the Regina audit.
The benchmark mode is separate so that correctness checks do not depend
on machine-specific timing thresholds.

To review the source delta at the pinned repository revision, use:

\begin{lstlisting}
git apply --check /path/to/affine_families_and_inheritance.patch
git apply /path/to/affine_families_and_inheritance.patch
\end{lstlisting}

The article and supporting research files can then be placed in the
repository's chosen report or incoming-document location. The distributed
patch edits the implementation paths; it does not choose or overwrite a
report number.

\subsection{Concrete example contracts}

The cyclic-family examples include a constrained annulus with optimum two,
a self-seam joining twelve locally disconnected annuli into a connected
member, an infeasible seam with a negative certificate, and a 10,589-bit
modulus whose stored hexadecimal witness verifies without sheet expansion.
The $K_4$ example keeps its designated-boundary requirements in a separate
metadata file, since that hard peripheral query is outside the optimizer's
strict input schema. Run:

\begin{lstlisting}
python scripts/check_family_examples.py
\end{lstlisting}

For direct package use, the public family functions are
\code{compile_family}, \code{optimize_cover_family},
\code{realize_family}, and \code{verify_cover_family_certificate} in
\code{fastunknot.cyclic_cover_family}. The command-line entry point accepts
one JSON input and emits the compiled family and arithmetic certificate.
The verifier ignores the producer's compilation and recomputes it from the
original geometry. It certifies feasibility and optimality, or infeasibility;
it does not certify the auxiliary assignment count or any topological
predicate outside the stated component objective.

The normal extractor accepts a JSON object with exactly
\code{triangulation} and \code{coordinates}. It emits the signed polygon
bands, prism bands, residual-cell incidence data, and provenance.
The optional \code{--orbit-input surface} and
\code{--orbit-input orientation} switches emit the corresponding interval
orbit problems. Those files are inputs to a subsequent orbit engine, not
claims that the engine has been run. The supplied normal examples include
the sharp two-tetrahedron fixture and a 128-tetrahedron Fibonacci meridian.

\subsection{Optional scalar controls and evidence boundaries}

Within the already optional Fitting backend, ordinary reuse defaults to
\code{fitting_reuse_endomorphisms=True}. It preserves the deterministic
candidate basis and search policy. The separately configurable primary,
locality, and one-generator paths default to false:

\begin{lstlisting}
from fastunknot.scalar_split import fitting_khovanov_rank
answer = fitting_khovanov_rank(
    pd,
    fitting_reuse_endomorphisms=True,
    fitting_primary=True,
    fitting_certify_local=True,
    fitting_monogenic_corners=True,
)
\end{lstlisting}

This snippet deliberately selects an experimental backend and its optional
mechanisms. It is not a recommendation to raise every size cap on arbitrary
inputs. Independent locality auditing uses
\code{verify_scalar_local_certificate}; unlike production inheritance,
that verifier recomputes the full commutant to check the dimension claim.
The different cost contracts of the arithmetic and scalar verifiers should
not be conflated.

\subsection{Provenance and continuing obligations}

The inherited polynomial-projector module, its fixtures and tests, and the
worker-communication correction are attributed by source path and hash.
The Fibonacci meridian family is likewise attributed to the earlier
compressed-certificate report. No external article PDF or downloaded slide
deck is bundled as a new research artifact. The new work follows the
repository's MIT No Attribution license; included pre-existing components
retain their accompanying license notices.

A claim ledger records proved statements, tested implementation contracts,
measured improvements, negative results, and unresolved global obligations.
The immediate deliverable is a set of exact reusable operations and an
integration-ready research record. A complete quasi-polynomial recognizer
still requires the explicit state-growth, representation-size, and
geometric-coverage arguments in Section~\ref{sec:global}.
